\chapter{Editing tracks}\label{ch:tracks}

The tracks are edited in the \texttt{TRACK EDIT} part (\key{F2}). The complete list of keys of the part is in
appendix~\ref{app:reference}; this chapter explains what they do.

\section{The cursor}

The arrow keys, \key{TAB} and \key{SHIFT+TAB}, \key{PAGE UP} and \key{PAGE DOWN} move the cursor in the tracks. The columns of each
track are, from the left: the note, the instrument, the volume and the speed. The key you type goes to the column under the cursor.

\begin{center}
\begin{tabular}{@{}lp{10.5cm}@{}}
\toprule
Column & What you type \\ \midrule
Note (\texttt{NNN}) & a note key; \key{numblock 1}--\key{6} change the octave of the note under the cursor; \key{BACKSPACE} deletes the note and the instrument \\
Instrument (\texttt{TT}) & \key{0}--\key{F}, from \texttt{\$00} to \texttt{\$3F}. If the note column of the line is empty the column behaves like a
  note column: you can type a note there \\
Volume (\texttt{vV}) & \key{0}--\key{F}. A volume can be written alone, with no note \\
Speed (\texttt{FSS}) & \key{0}--\key{F}: a new speed for the song from this line on, from \texttt{\$01} to \texttt{\$FF} \\
\bottomrule
\end{tabular}
\end{center}

If the \emph{respect volume} mode (\key{F11}) is on, a note typed on a line that has a volume does not replace the volume.

\section{The length of a track}

A track is as long as its \emph{end line}: \key{CONTROL+END} sets or clears it. \key{HOME} and \key{END} go to the start and
the end of the track; \key{CONTROL+HOME} sets the start of a \emph{wise loop}, the line the track goes back to when it
reaches its end, so that a repeated phrase is written once. The header of the track shows both numbers
(chapter~\ref{ch:window}).

\key{INSERT} and \key{DELETE} (or \key{CONTROL+I} and \key{CONTROL+U}) insert and delete lines in the track.

\section{Playing while editing}

\begin{itemize}
  \item \key{ENTER} plays the note at the cursor, \key{SHIFT+ENTER} plays it and makes its instrument and volume the active ones
        (a quick way to ``pick up'' a sound from a song), and \key{CONTROL+ENTER} plays all the notes of the line.
  \item \key{SHIFT+tonekey} plays a note with the active instrument and volume on the current channel, without writing it, and
        \key{SHIFT+CONTROL+tonekey} plays it on both stereo channels.
  \item \key{F9} switches the current channel off and on, \key{CONTROL+F9} plays it alone, \key{CONTROL+1}--\key{8} toggle
        the channels one by one.
\end{itemize}

\section{Blocks}

A \emph{block} is a selection of lines in a track.

\begin{enumerate}
  \item Select with \key{SHIFT+UP}, \key{SHIFT+DOWN}, \key{SHIFT+HOME} and \key{SHIFT+END}; \key{CONTROL+A} selects all the data of the
        track. \key{ESC} deselects.
  \item Copy (\key{CONTROL+C}), cut (\key{CONTROL+X}), paste (\key{CONTROL+V}), merge with the content under the cursor
        (\key{CONTROL+M}) and exchange the block with the clipboard (\key{CONTROL+E}). With no block the data at the cursor
        is taken.
  \item \key{CONTROL+B} restores the block from a backup that is made as soon as the block is started: the way back from a wrong
        modification.
\end{enumerate}

The blocks can also be changed with keys, either on every line of the block, or only on the lines with the active instrument
(\key{SHIFT+CONTROL+A} chooses): \key{CONTROL+F1}/\key{F2} transpose by semitones, \key{CONTROL+F3}/\key{F4} by octaves,
\key{SHIFT+CONTROL+LEFT}/\key{RIGHT} change the instrument and \key{SHIFT+CONTROL+UP}/\key{DOWN} the volume. The block toolbar has the same
buttons.

\subsection{The Effects and Tools dialog}

\key{CONTROL+F} (\menu{Block $\rightarrow$ Effects/Tools\ldots}) opens the window of figure~\ref{fig:dialog-block-effects}: a
choice of effects for all the data of the block, with the fields of the chosen effect. The effects include fading the volume in
or out between two levels in percent; \menu{Try} applies the effect, \menu{Restore} takes it back, \menu{Play/Stop} plays the
block while you are adjusting it and \menu{Default} sets the fields to their default values.

\dialog{dialog-block-effects}{The block effects (\key{CONTROL+F}).}{7cm}

\section{Tracks as objects}

The \menu{Track} menu works on a whole track: copy, paste, cut, delete, load a track from a file or save it as a
text file, ask for information about the current track, \emph{search and build a wise loop} in a track (the program finds the
repeats and writes them once) or \emph{expand} it again, and the clean-up operations that remove duplicated or unused
tracks from the song. The operations on all the tracks change the whole song; save it before using them.

\section{Undo}

\key{CONTROL+Z} undoes the last change and \key{CONTROL+Y} does it again. \menu{Edit $\rightarrow$ Clear Undo \& Redo History}
forgets the changes made so far.
